%%%%%%%%%%%%Outline%%%%%%%%%%%%%%%%%%

%%% Section 8 Related Work

%%%
% message<eBPF bytecode optimization: bounded by existing ISA>
% bytecode optimizers improve eBPF programs without changing the ISA
% K2: stochastic search for shorter bytecode sequences
% Merlin: expert-designed rewrite rules at LLVM-IR and bytecode levels
% EPSO: offline enumerative superoptimization with cached rules
% ePass: SSA-based compiler framework cooperating with verifier at load time
% KFuse: post-verification merging of tail-call chains
% all operate within existing eBPF ISA
% \lang lifts that constraint by adding opcodes the eBPF JIT compiler maps to native instructions

%%%
% message<extending or relocating the eBPF pipeline>
% several systems modify the eBPF pipeline itself
% bpftime: userspace, near-native performance
% hXDP: FPGA ISA extension, no verifier-checked semantic equivalence
% Craun 2023: move verification off embedded devices to separate host
% contrast: \lang's decoupling is representational, not spatial
% standardizing new opcodes through existing path is costly
% Wasm SIMD: years of multi-engine standardization; eBPF ISA revisions: kernel/toolchain changes, CVEs
% \lang avoids this cost: proof sequences let verifier check new opcodes without core changes

%%%
% message<verified compilation and proof-carrying code>
% verified compilation and proof-carrying code provide correctness techniques \lang uses
% CompCert: simulation-based refinement methodology
% Alive2: per-transformation refinement for LLVM
% Jitterbug: per-instruction eBPF JIT refinement via SMT
% \lang per-opcode refinement check: same per-instruction structure, module-defined opcodes not fixed ISA
% Agni: verifier range-analysis soundness, prerequisite for \lang safety
% BCF carries abstraction-refinement proofs, \lang carries semantic equivalence proofs
% Necula PCC: safely load untrusted code into kernel, direct ancestor
% \lang reuses existing eBPF verifier instead of deploying a new checker

%%%%%%%%%%%%Outline%%%%%%%%%%%%%%%%%%

\section{Related Work}
\label{sec:related}

We compare \lang to prior work along three axes: optimizing within the existing eBPF ISA, extending or relocating the eBPF pipeline, and formally verified compilation.

\noindent\textbf{eBPF bytecode optimization.}
Bytecode optimizers improve eBPF programs without changing the ISA. K2~\cite{xu2021synthesizing} uses stochastic search to find shorter equivalent bytecode sequences. Merlin~\cite{mao2024merlin} applies expert-designed rewrite rules at the LLVM-IR and bytecode levels. EPSO~\cite{zhu2025epso} uses offline enumerative superoptimization with cached rules. ePass~\cite{xiang2025epass} adds an SSA-based compiler framework that cooperates with the eBPF verifier at load time. KFuse~\cite{kfuse} merges tail-call chains into a single program after verification. These optimizers all operate within the existing eBPF ISA. \lang lifts that constraint by adding opcodes that the eBPF JIT compiler maps directly to native instructions.

\noindent\textbf{Extending or relocating the eBPF pipeline.}
Several systems modify the eBPF pipeline itself rather than optimizing programs within the existing ISA. bpftime~\cite{zheng2025extending} runs eBPF programs in userspace, achieving near-native performance. hXDP~\cite{brunella2020hxdp} extends the eBPF ISA for FPGA targets, without verifier-checked semantic equivalence. Craun et al.~\cite{craun2023enabling} move eBPF verification off embedded devices to a separate host. Their spatial decoupling contrasts with \lang's representational decoupling: verification and execution consume different bytecode representations on the same kernel. Standardizing new bytecode opcodes through the existing path is costly. WebAssembly's SIMD extension~\cite{haas2017bringing} required years of multi-engine standardization, and eBPF's own ISA revisions have required substantial kernel and toolchain changes and introduced CVEs. \lang avoids this cost: proof sequences let the eBPF verifier check new opcodes with no core changes, and adding a new instruction requires only a kernel module.

\noindent\textbf{Bytecode design in mainstream runtimes.}
Mainstream language runtimes balance verification strictness against JIT optimization scope differently. LLVM~\cite{lattner2004llvm} performs no safety verification on its IR and runs a full middle-end and backend optimization pipeline. The JVM relies on a load-time bytecode verifier for type safety, and HotSpot~\cite{paleczny2001java} compiles verified bytecode through a tiered JIT with profile-guided optimization. Wasm~\cite{haas2017bringing} validates control flow and operand types at load time, and Wasm engines perform cross-instruction JIT optimization. Among these runtimes, only eBPF combines a strict per-opcode abstract interpretation (tracking types, value ranges, and memory provenance) with a JIT limited to encoding-level optimizations. \lang adds a representational split that none of these runtimes currently uses: the verifier continues per-opcode abstract interpretation on lowered bytecode, while the eBPF JIT compiler emits hardware-specialized native code through module-defined opcodes.

\noindent\textbf{Verified compilation and proof-carrying code.}
Verified compilation and proof-carrying code provide the correctness techniques \lang uses. CompCert~\cite{leroy2009formal} established the simulation-based refinement methodology that \lang instantiates per opcode. Alive2~\cite{lopes2021alive2} performs per-transformation refinement checking for LLVM. Jitterbug~\cite{nelson2020specification} verifies per-instruction eBPF JIT translation via SMT-based refinement. \lang's per-opcode refinement check has the same per-instruction structure but targets module-defined opcodes rather than the fixed ISA. Agni~\cite{vishwanathan2023verifying} verifies the eBPF verifier's range analysis, whose soundness is a prerequisite for \lang's safety argument. BCF~\cite{sun2025prove} carries abstraction-refinement proofs to improve verifier precision, while \lang carries semantic equivalence proofs between the lowered and native representations. Necula's proof-carrying code~\cite{necula1997pcc} is \lang's most direct ancestor: both safely load untrusted code into the kernel. PCC requires deploying a new proof checker, while \lang reuses the existing eBPF verifier.